% ED26-SRC-036-014, Davis 2008 Physical Reasoning; assigned coordination filename.
\chapter{Physical Action, Kinematics and Liquids}
\label{sec:p16-kr2008-action-and-causation}

Ernest Davis's physical-reasoning chapter concerns the laws and representations
needed to understand actions in a geometrically and temporally continuous
world \cite[pp.~597--620]{Davis2008PhysicalReasoning}.
An action language can describe a change without explaining why an actual
manipulator can perform it. The STRIPS blocks world does not even specify
the swept positions of a block while it moves. Conversely, an accurate
scientific simulator can be too specialized or detailed for a commonsense
question. Davis's concern is a physical ontology and usable domain theory,
not a replacement definition of nonmonotonic causal explanation or C+.

\section{Physical assumptions, fluents and architectures}
\label{sec:p16-phys-architectures}

The chapter distinguishes significant-event/region reasoning from numerical
discretization, and asks for inference in several directions, explanation,
action, continuous geometry and interaction between models. It separates
two defeasible commitments: closure of the relevant entities and mechanisms,
and adequacy of an idealization. Knowing all objects does not justify making
them rigid; making them rigid does not establish that all relevant forces
are known. A qualitative answer can still depend on precise numerical data,
as in the car stopping just short of a cliff
\cite[pp.~597--600]{Davis2008PhysicalReasoning}.

A Davis fluent is a time-varying entity. \(Temperature(o)\) and
\(Place(o)\) are functions of time, while \(On(a,b)\) is Boolean-valued.
A parameter is a numeric-valued fluent. Addition and differentiation act on
the time functions, not on an isolated qualitative label.

In the component architecture, a component has ports with parameters and
algebraic/differential constraints. Nodes group connected ports and add
domain-specific connection constraints. The component/connection structure
is fixed; the parameters vary. For an electrical node, incoming currents
sum to zero and port voltages agree. An agent can be an exogenous component
which chooses an input waveform. Static, initial-value, signal-response and
comparative problems use the same physical constraints differently
\cite[pp.~600--601]{Davis2008PhysicalReasoning}.

\begin{hbdefinition}{def:p16-phys-rc}{An actual construction from the source's port equations}
The source resistor has \(InCurrent(a)=-InCurrent(b)\) and
\(Voltage(b)-Voltage(a)=R\,InCurrent(b)\). Its capacitor has the same
current balance and \(InCurrent(b)=C(Voltage(b)-Voltage(a))'\).
Here \(R,C>0\) are constant. Connect a resistor from a chosen voltage
\(u(t)\) to a capacitor node of voltage \(v(t)\), whose other port is
at zero voltage. Orient the resistor's \(b\) port at \(u\) and its
\(a\) port at \(v\). At their shared node, Kirchhoff current balance
equates the resistor current \((u-v)/R\) with \(Cv'\).
This ideal RC construction assumes no other branch current, sources or
changing component characteristics.
\end{hbdefinition}

\begin{hbtheorem}{thm:p16-phys-rc}{proved=annex:p16.phys-rc}{Composed signal response with all boundary data stated}
For continuous \(u\) on \([0,T]\), \(T>0\), and specified \(v(0)=v_0\), the
port/node constraints give \(RCv'+v=u\) and their unique differentiable
solution is
\[
 v(t)=e^{-t/(RC)}v_0+
       \frac1{RC}\int_0^t e^{-(t-s)/(RC)}u(s)\,ds.
\]
An unspecified input waveform or initial capacitor voltage does not determine
a unique response. The construction concerns this physical idealization;
it is not a proof of completeness of a qualitative sign predictor.
\end{hbtheorem}

In the process architecture, a process lasts while its conditions hold
and contributes to a derivative. An optional action takes place at an
instant if feasible; a natural event must occur when its triggering conditions
hold. Indirect parameter relations are distinct from direct influences.
Davis retains the no-direct-and-indirect-overlap and acyclic-indirect-dependency
restrictions of the process framework, already developed in the separate
qualitative-process treatment \cite[pp.~601--602]{Davis2008PhysicalReasoning}.

\begin{hbexample}{def:p16-phys-bucket}{The source's process/action/event distinction}
A tap fills a bucket while open, positioned above it and while the bucket
is not full. Other taps contribute positively and a leak negatively to
the volume derivative. Turning on the tap is an agent's optional action,
requiring proximity and a closed tap; it changes the tap's discrete state.
The weak bucket's bottom falling out at half capacity is instead a mandatory
event, changing the bucket into a cylinder and a pan.

For a stipulated constant positive net rate \(r\), initial volume
\(V_0<B/2\), and retained conditions, the pre-event history is
\(V(t)=V_0+rt\); equality with half capacity occurs at
\((B/2-V_0)/r\). Choosing to close the tap before this time changes the
active process set; reaching the trigger is not a choice to postpone the
event. Simultaneous interventions and event priority need separate semantics.
The event also changes the object structure, not merely a value on a fixed
component list. The source's envisionment can admit paths without physical
realizations; it does not identify every graph path with an actual history.
\end{hbexample}

\section{Rigid shape, pose, place and configuration space}
\label{sec:p16-phys-kinematics}

The source kinematic theory preserves shape, requires continuous position
and prevents interpenetration. It does not predict that a book stays on
a shelf rather than floating away. Shape is the occupied region in a
reference position; position is the rigid mapping from that frame; place
is the resulting occupied region. Figure~14.1 distinguishes those objects
\cite[pp.~603--605]{Davis2008PhysicalReasoning}.

\begin{hbdefinition}{def:p16-phys-configuration}{A precise rigid-body configuration model}
For the following constructions, take finitely many objects with nonempty
compact, connected, regular closed shapes \(S_o\subseteq\mathbb R^d\),
\(d=2\) or \(3\). Regular closed means closure of the interior.
A pose is \((Q_o,b_o)\), where \(Q_o^TQ_o=I\), \(\det Q_o=1\),
and \(b_o\in\mathbb R^d\). Proper rotations exclude reflections.
The source functions have the explicit evaluation
\[
 Shape(o)=S_o,\quad Position(o)(t)=(Q_o(t),b_o(t)),\quad
 Place(o)(t)=Q_o(t)S_o+b_o(t).
\]
A configuration is the tuple of poses. Its feasible set \(\mathcal F\)
requires disjoint \emph{interiors} of occupied regions for distinct objects;
contact at boundaries is allowed. Additional genuine joint constraints,
if any, must be declared. Give pose tuples their ordinary matrix/vector
topology, not discontinuous Euler-angle coordinates.
\end{hbdefinition}

Compactness is supplied here for the occupied-set continuity argument;
it is not silently inferred from regular closedness. The source's rotation
description, rather than arbitrary orthogonal reflections, motivates the
proper-rotation convention. Likewise forbidding any intersection of closed
regions would incorrectly forbid the contact needed by mechanics.

\begin{hbtheorem}{thm:p16-phys-kinematics}{proved=annex:p16.phys-kinematics}{Rigid invariants and the exact attainable-history construction}
Each pose preserves compactness, connectedness and regular closedness of a
shape. A continuous pose history gives continuous occupied regions in
Hausdorff distance. In this stated kinematic model, a configuration \(q_1\)
is attainable from \(q_0\) exactly when a continuous path in
\(\mathcal F\) joins them. Every such path constructs a full model
history via the source's shape/position/place functions; every model history
gives such a path. This is a path component, not automatically an ordinary
connected component for arbitrary shape descriptions.
\end{hbtheorem}

The source's computational organization is to obtain the feasible
configuration space from exact shape descriptions, select the initial
attainable component, divide it into physically significant regions and
analyze continuity/connectivity. Low-dimensional mechanisms make this useful;
a few moving parts can nevertheless yield a difficult global geometry.
No finite sign table or mere count of degrees of freedom proves that this
construction terminates for arbitrary shapes.

\section{Qualitative cells and physically feasible manipulator paths}
\label{sec:p16-phys-cells-actions}

\begin{hbtheorem}{thm:p16-phys-cells}{proved=annex:p16.phys-cells}{An exact finite cell-cover construction, with its premises}
Suppose \(\mathcal F\) has a finite cover by nonempty path-connected
subsets \(C_1,\ldots,C_n\), each closed relative to \(\mathcal F\).
Construct a graph with an edge when \(C_i\cap C_j\ne\varnothing\).
Two configurations are path-reachable exactly when cells containing them
are connected in this graph. A graph path constructs a continuous physical
path by joining intersection witnesses within each cell.

If qualitative cells are disconnected, or edges only mean that closures
touch at a point excluded from \(\mathcal F\), the conclusion can fail.
Computing this exact cover and its intersections remains a geometric task.
\end{hbtheorem}

A qualitative category can accidentally unite two disconnected corridors.
An edge entering one corridor and another edge leaving the other then gives
a false qualitative route. Contact with an excluded collision configuration
also does not establish a usable passage. These are why partition labels
alone cannot establish global physical realizability.

\begin{hbtheorem}{thm:p16-phys-action-lift}{proved=annex:p16.phys-action-lift}{The source action-feasibility requirement as a continuous lift}
Fix the initial configuration and designate a subset of objects manipulable.
A specified continuous pose motion \(m(t)\) for them is kinematically
feasible exactly when it has a continuous lift \(q(t)\in\mathcal F\)
from the initial configuration whose manipulable coordinates equal \(m(t)\).
Feasibility of every individual frame does not suffice.

For an explicit abstract constrained mechanism,
\(\mathcal F=\{(x,0):x\le0\}\cup\{(x,1):x\ge0\}\), every
specified \(x(t)=2t-1\), \(0\le t\le1\), has a feasible frame.
But no continuous lift starts at \((-1,0)\) and performs the whole motion.
The constraints, not a statistical intervention variable, define this example.
\end{hbtheorem}

Davis's projection problem specifies shapes, initial positions and motions
of manipulable objects. Unmanipulated objects must accommodate those motions
along a feasible path, if there is one. This does not choose their unique
trajectory or explain forces. Figure~14.2's pinned gears are another
important abstraction: remove pin/frame geometry and impose fixed centers.
That constraint has to be justified by the concrete joint and planar-motion
assumptions; permitting a three-dimensional lift can invalidate it.

\section{Dynamics, contact history and the limits of a physical explanation}
\label{sec:p16-phys-dynamics}

Rigid dynamics adds Newton's force/momentum laws, contact normals, static
and sliding friction, impact assumptions, gravity, fixed objects and chosen
manipulator forces. Force closure is an assumption. Kinematics alone permits
a book to leave its shelf along a collision-free upward path; statics can
instead establish equilibrium. Davis distinguishes point-object prediction,
statics and quasi-statics, rather than calling all of them rigid dynamics
\cite[pp.~605--608]{Davis2008PhysicalReasoning}.

In quasi-statics momentum is neglected relative to driving and dissipative
forces. This can support sensorless pushing or qualitative mechanical
analysis, but it is not arbitrary slow-time simulation. The simplification
``objects move only when necessary'' needs the relevant contact, force and
kinematic restrictions. Prescribing a force, prescribing a motion and using
feedback-dependent compliance are different descriptions of manipulation.
The chapter does not supply a universal high-level language for them or an
unconditional existence/uniqueness theorem for all rigid-body impacts.
Its discussion of unsettled collision laws is part of the 2008 source's
assessment, not a claim about today's state of mechanics.

\begin{hbdefinition}{def:p16-phys-nail}{A contact-force model for Figure 14.3}
Idealize the vertical nail by two symmetric, equal-height side contacts.
Each supplies a horizontal normal force of specified magnitude \(N\ge0\).
The forces oppose and their torques cancel. Upward static friction at each
contact is between zero and \(\mu N\), where \(\mu\ge0\).
Gravity acts through the central mass, with weight \(mg>0\); no other
forces act. This is a stated static contact model, not a prediction of
\(N\) from position and velocity alone.
\end{hbdefinition}

\begin{hbtheorem}{thm:p16-phys-nail}{proved=annex:p16.phys-nail}{Same geometry, different force boundary conditions}
A static equilibrium of this model exists exactly when \(mg\le2\mu N\).
Choosing equal friction forces \(mg/2\) balances both vertical force and
torque when this inequality holds. With zero normal force no such equilibrium
exists; when \(\mu>0\), sufficiently preloaded contacts can hold the same
positioned nail.
Thus its geometry and zero initial velocity do not specify this static outcome.
\end{hbtheorem}

Davis's example distinguishes a nail placed in a hole from one driven into
it. The latter has substantial normal/frictional force information not
encoded by its visible pose. The static construction makes that omitted
boundary datum explicit; it does not infer every preload or collision force.

The source's funnel argument excludes upward escape by energy, a stable
interior rest by geometry, and perpetual internal motion by dissipation.
Those are three separate obligations. The marble must fit the outlet;
adhesion, jamming, other supports or an external drive change the problem.
Moreover, energy decreasing by itself does not establish finite exit.
The following supplies a precise sufficient progress condition rather than
claiming that the source's short prose is a complete general dynamics proof
\cite[pp.~608,615, referral~9]{Davis2008PhysicalReasoning}.

\begin{hbtheorem}{thm:p16-phys-exit}{proved=annex:p16.phys-exit}{A full bounded energy-and-progress argument}
Let a jointly continuous forward flow \(\Phi_t\) describe a metric
phase-space model, with \(\Phi_0\) the identity and
\(\Phi_{t+s}=\Phi_t\circ\Phi_s\) for nonnegative times,
and let \(K\) be a compact candidate trap. Let a continuous energy \(E\)
have a continuous derivative along the flow, \(D\le0\). If there is
no nonempty forward-invariant subset of \(K\cap\{D=0\}\), no trajectory
can remain in \(K\) forever. If an exit face has potential energy strictly
greater than initial total energy, with total energy at least that potential,
the trajectory cannot exit through that face.

For the smooth damped model \(q'=v,\ v'=-\nabla U(q)-\gamma v\),
\(\gamma>0\), energy \(\|v\|^2/2+U(q)\) has
\(D=-\gamma\|v\|^2\). If \(K\) contains no zero-velocity
equilibrium \(\nabla U(q)=0\), it satisfies the progress condition.
This is a continuous-flow result, not a theorem covering arbitrary nonsmooth
impact laws. The abstract exit system \(x'=-x\), \(0<x\le1\),
with exit at zero and energy \(x^2\), decreases energy and has no interior
stationary point yet does not exit at finite time.
\end{hbtheorem}

Pseudo-objects in Davis's logical work are geometric entities carried with
an object, such as a doughnut's hole or a center of mass. They permit
statements about a passage or material support which a bare object name
does not express. A hole is not another solid to which noninterpenetration
is applied. The complete funnel axiomatization referred to by the
chapter is a distinct original work; the bounded progress proof above does
not silently close it.

\section{Liquids: regions, chunks, particles and moving control boundaries}
\label{sec:p16-phys-liquid-ontology}

Liquids lack a fixed rigid shape and are not naturally individuated into
permanent object-sized bodies. Davis contrasts three complementary
representations. Region quantities use \(Volume(l,r)\) and oriented
\(Flow(l,b)\). Material chunks use \(Place(c)\), which can be disconnected.
Particle/molecular descriptions follow many small units and need an interaction
model producing liquid behavior. A material reference-time description uses
\(LiquidTrajectory(X,L)\) for the later position of the liquid particle
at \(X\) initially. A material chunk and a spatial observation region
are not the same kind of entity
\cite[pp.~608--610]{Davis2008PhysicalReasoning}.

The source's liquid kinematics assumes no mixing/phase change,
incompressibility, continuous motion and noninterpenetration with solids
or other liquids. It preserves each material quantity's volume. A moving
container can be used as the region argument, but the boundary frame must
then enter the flux definition.

\begin{hbdefinition}{def:p16-phys-liquid-balance}{A smooth material/control-region balance with explicit frame}
For this analytic portion take compact liquid and control regions \(L(t)\)
and \(r(t)\), with compact twice-continuously-differentiable boundaries
meeting transversely. Assume their
motion is differentiable, with liquid boundary velocity \(v\), control
boundary velocity \(w\), and a continuously differentiable liquid velocity
field extending to its boundary with \(\nabla\cdot v=0\).
Normals on \(\partial r\) point outward. Liquid volume in the control
region is \(V(t)=\operatorname{vol}(L(t)\cap r(t))\); its relative
outflow is
\[
 F_{rel}(t)=\int_{\partial r(t)\cap L(t)}(v-w)\cdot n_r\,dA.
\]
Smoothness here includes the local boundary-layer differentiability used
in the annex; intersections have codimension two, not coincident faces.
Contact/topology-change cases require their appropriate piecewise treatment.
\end{hbdefinition}

\begin{hbtheorem}{thm:p16-phys-liquid-balance}{proved=annex:p16.phys-liquid-balance}{Relative outflow and material conservation}
At that scope \(V'=-F_{rel}\); reversing orientation reverses outflow.
For a fixed control region \(w=0\), this is the laboratory-frame balance.
A material liquid region preserves volume. For a moving control region,
using laboratory flux without boundary motion can give a false derivative.
For example, in a unit-cross-section one-dimensional model, stationary
liquid \([0,10]\) and control interval \([t-1,t+1]\), \(-1<t<1\),
have liquid volume \(t+1\), although laboratory liquid velocity is zero.
\end{hbtheorem}

The source's \(Volume\)/\(Flow\) equation therefore needs ``Flow''
to mean outward flux \emph{relative to the specified moving surface},
or its region must be fixed. Conservation of a chunk's total volume does
not make its volume in every spatial region constant. Nor does a region-
flow representation by itself determine a unique fluid dynamics.

\section{Hydrostatic levels, capacity and qualitative liquid histories}
\label{sec:p16-phys-liquid-statics}

Ignore wetting and surface tension and take one connected body of liquid
at rest under uniform gravity, with constant density \(\rho>0\).
Free surfaces are horizontal. Common open air makes their heights equal;
different trapped-gas pressures can make them unequal. Figure~14.4's
surfaces A and B meet the atmosphere, while C is pressed down by enclosed
gas. Omitting that gas loses a necessary physical entity.

\begin{hbtheorem}{thm:p16-phys-liquid-level}{proved=annex:p16.phys-liquid-level}{The actual pressure/height and capacity conditions}
On a connected open fluid interior, assume pressure is continuously
differentiable inside and continuous up to the relevant free surfaces,
and assume the hydrostatic law
\(\nabla p=-\rho g\,e_z\), \(g>0\). Surfaces at heights
\(h_i,h_j\) with gas pressures \(p_i,p_j\) satisfy
\[
 h_i-h_j=\frac{p_j-p_i}{\rho g}.
\]
Equal gas pressures give equal heights; higher pressure lowers its surface.
For an open watertight container with no trapped gas or inaccessible cavity,
let \(A(h)=\int_{h_0}^{h}a(z)\,dz\) be its below-level capacity,
where \(h_0<h_r\) are finite and \(a\ge0\) is integrable and has positive integral on every
nonempty height interval up to the lowest opening \(h_r\).
Every \(0\le V\le A(h_r)\) has exactly one level \(A(h)=V\).
If \(V>A(h_r)\), no wholly contained hydrostatic state satisfies these
premises. Applying this to slow tilting still needs the quasi-static
assumption; inertia, sloshing or capillarity can invalidate that application.
\end{hbtheorem}

The source's tilted-cup rule is a capacity comparison, unlike the geometry
needed to tip a rigid object from a box. A slowly moving piston can sometimes
preserve hydrostatic constraints as liquid/gas moves; spilling is a case
where that simple quasi-static maintenance becomes impossible. Kim's liquid/
mechanical model also represents one-way-valve motion caused by liquid
\cite[pp.~608--611]{Davis2008PhysicalReasoning}.

Hayes' history ontology partitions liquid quantities and their space--time
histories using spatial distribution and motion categories. Davis's table
has fifteen nonempty categories. Its examples and question marks are retained
below; blank cells are not a proof that all unlisted physical phenomena
are impossible.

\begin{table}[ht]
\centering
\begin{tabular}{p{.19\linewidth}p{.22\linewidth}p{.22\linewidth}p{.22\linewidth}}
 & Lazy still & Lazy moving & Energetic moving\\\hline
Bulk on surface & Wet surface & Flow down a surface, e.g. sloping roof
 & Waves lapping a shore (?); jet hitting a surface (?)\\
Bulk in space & Contained in a container & Flow along a channel, e.g. river
 & Pumped along a pipeline\\
Bulk unsupported & -- & Falling column, e.g. pouring from a jug
 & Waterspout, fountain, hose jet\\
Divided on surface & Dew, drops on a surface & -- & --\\
Divided in space & Mist filling a valley & Mist rolling down a valley
 & Steam or mist blown along a tube\\
Divided unsupported & Mist, cloud & Rain, shower & Spray, splash, driving rain
\end{tabular}
\caption{Hayes' liquid-history categories as surveyed in Davis's Table 14.1.}
\label{tab:p16-phys-liquid-histories}
\end{table}

A falling history has downward outflow at its bottom face; a pool's faces
other than its upper horizontal one meet solid histories. Those spatial
and temporal face relations are representation obligations, not just state
names. Davis explicitly says that the proposed liquid-history axiomatization
was not completed and that its intended spatial language was insufficient.
No complete liquid simulator or exhaustive formal classification theorem
is inferred from the table \cite[pp.~610--611]{Davis2008PhysicalReasoning}.

\section{Abstraction is a physical commitment, not a guarantee of easier reasoning}
\label{sec:p16-phys-abstraction}

Davis separates several relationships between physical theories: added
constraints, added entities, limiting cases, greater numerical precision,
discretization into selected states/events, smoothing discrete phenomena,
and radically different conceptualizations such as wave/particle accounts.
These are not all the same logical weakening. Ignoring thermal expansion
or approximating a ball by a point can remove a decisive distinction.
Conversely a stronger static or frictional theory can simplify the actual
question by proving that nothing moves
\cite[pp.~611--614]{Davis2008PhysicalReasoning}.

\begin{hbtheorem}{thm:p16-phys-abstraction}{proved=annex:p16.phys-abstraction}{What weakening a theory or shrinking an object preserves}
For same-vocabulary theories \(U=V\cup A\), every \(U\)-model is
a \(V\)-model; \(V\)-consequences remain \(U\)-consequences.
Consistency with \(V\) need not imply consistency with \(U\).
This ordering alone makes no assertion about computational difficulty or
physical validity of the added constraints.

For a disk of radius \(r>0\) in a planar straight corridor of width
\(w>0\), between fixed parallel walls, its center can be feasible inside
the corridor exactly when \(w\ge2r\). Replacing the disk by a point
admits a corridor path for every \(w>0\), including \(w<2r\) where
the original disk has no feasible center. At equality contact is allowed
kinematically; friction or force limits remain different questions.
\end{hbtheorem}

The source's loosely meshed gear train has two coordinated paths, for
opposite driving directions, separated by a narrow tube of free-play
configurations. Collapsing that tube to a central path can be useful for
one task and wrong when reversal/backlash matters. Treating a table as one
object is valid only while its internal joints remain fixed; it fails when
the table breaks. A hierarchy can also share subparts, defeating an
assumption that high-level components are independent.

The pinned-gear abstraction is computationally straightforward when all
later reasoning uses its computed constrained configuration space. In a
logical theory, the generated constraint must either be reified with a
semantics or added as an axiom justified from the concrete system.
Combining it with the original geometry remains a consistency obligation.
It cannot simply be asserted because the downstream solver handles it.

\begin{hbdefinition}{def:p16-phys-braking}{A quantitative margin for the source's qualitative cliff question}
For this worked idealization, treat the car's relevant front position as
a point initially zero, moving toward an edge at distance \(L>0\).
Braking gives constant deceleration \(a>0\) until speed reaches zero,
then holds it stationary, with no new driving force. Initial speed
\(v_0\ge0\) and deceleration are independently permitted to vary over
\([v_{min},v_{max}]\) and \([a_{min},a_{max}]\), with
\(0\le v_{min}\le v_{max}<\infty\) and
\(0<a_{min}\le a_{max}<\infty\). This excludes reaction delay, slope variation, finite
body geometry and changes of traction. Those are not inferred negligible.
\end{hbdefinition}

\begin{hbtheorem}{thm:p16-phys-braking}{proved=annex:p16.phys-braking}{Exactly when the qualitative answer is determined by those bounds}
The stopping distance is \(d=v_0^2/(2a)\), and the full range over
the specified parameter rectangle is
\[
 d_{min}=\frac{v_{min}^2}{2a_{max}},\qquad
 d_{max}=\frac{v_{max}^2}{2a_{min}}.
\]
Stopping strictly before the edge is guaranteed iff \(d_{max}<L\),
and possible iff \(d_{min}<L\). Passing the edge is guaranteed iff
\(d_{min}>L\), and possible iff \(d_{max}>L\). If
\(d_{min}<L<d_{max}\), both outcomes are realizable. Equality is
a boundary case, not a proof of either strict outcome. Correlations among
the parameters can invalidate these rectangle extrema.
\end{hbtheorem}

This is why a qualitative output does not eliminate the need for numerical
information near a threshold. It is an analysis of the declared model,
not a driving recommendation or a prediction that a real car will crash.

\section{The source's historical and representational comparisons}
\label{sec:p16-phys-history}

Davis's historical landmarks are the qualitative-program literature,
Hayes' Naive Physics Manifesto/liquid ontology, and the 1987 configuration-
space work of Faltings and Joskowicz. The qualitative-calculus development
belongs to Chapter 9; its original and reprinted source versions are not
silently merged with later QSIM accounts.

The logic-based branch includes Schmolze's actions/events/processes/containers
and liquid granules; Sandewall's combination of nonmonotonic logic with
continuous motion; the distinct egg-cracking axiomatizations of Lifschitz,
Morgenstern and Shanahan; and Bennett and colleagues' geometric treatment of
rigid motion. Davis's cutting work distinguishes an object which changes
shape and splits from material chunks which persist until penetrated by the
blade. These are different persistence/individuation choices, not variant
spellings of one object ontology. Hayes' cluster programme and Davis's later
microworld methodology are also distinct
\cite[pp.~614--615]{Davis2008PhysicalReasoning}.

The mechanics referrals distinguish exact geometric simulation from qualitative
configuration analysis. CLOCK constructs an exact configuration space before
abstraction and force-based qualitative prediction. Stahovich, Davis and
Shrobe use planar one-degree-of-freedom objects, negligible inertia and
inelastic collisions; the input qc-space contains pairwise feasible/contact
regions. Davis explicitly notes that its proposed extraction from an informal
sketch is not specified clearly. Their simulation is not a complete algorithm
for arbitrary three-dimensional rigid bodies or arbitrary input images
\cite[pp.~615--616]{Davis2008PhysicalReasoning}.

The multiple-model referrals address reformulation and approximation.
The other-domain survey includes particle/hybrid liquid models, qualitative
magnetic-flux reasoning and narrowly specialized scientific systems.
Its Project Halo discussion distinguishes held-out answer/explanation
evaluation from evidence of general representational coverage: a timeless
chemical equation need not represent the temporal replacement of reactants
by products. Davis's criticism is a 2008 source assessment, not an independent
current evaluation of those systems. Bobrow's reprint volume, Weld/de Kleer's
readings and Kuipers' 1994 book are distinct source works
\cite[pp.~616--620]{Davis2008PhysicalReasoning}.

\begin{hbmethod}{def:p16-phys-consultation}{A physical-action modeling specification}
Identify the actual bodies, material quantities, boundaries and agents.
State closure and idealization separately. Decide whether the task needs
kinematics, equilibrium, quasi-statics or a full dynamic model. For motion,
construct the relevant configuration constraints and a continuous path, not
just endpoint feasibility. For dynamics, record contact/preload, force and
impact assumptions. For liquids, distinguish material chunks from moving
control regions and use the proper flux frame and gas/pressure information.
When changing abstraction, identify the discarded distinction and check its
effect on the query; do not equate a quotient graph or state classification
with all physically realizable trajectories.
\end{hbmethod}

\index{physical reasoning}\index{configuration space}\index{kinematics}
\index{liquid!ontology}\index{liquid!moving boundary}\index{abstraction!physical}
